\section{动机实验的协议与对照}
\label{app:protocol}

\paragraph{共同设置与数据用途。}
所有动机实验使用 CLIP ViT-B/32~\citep{radford2021learning}、八个固定任务专家（Cars、DTD、EuroSAT、GTSRB、MNIST、RESISC45、SVHN、SUN397）及 seed 0。每任务统计集包含 256 张图像，另外 256 张不重叠图像用于式~\eqref{eq:behavior} 的选步与选轮；每轮 CG100，最多四轮。所有动机面板均使用卡方归一化后的 KL 修正统计。完整测试只用于报告，不参与更新选择。
多起点研究在 validation 的前 256 张图像上诊断，并在其中 64 张记录 block 表征；多轮对照另用固定随机选取的 256 张 validation 图像诊断。因此，两组诊断指标仅在各自配对协议内比较。

\paragraph{图 1c 的“度量变化”具体是什么？}
对层 $\ell$、任务 $t$ 和当前状态 $r$，记按求解规则处理后的双侧度量为
$M_t^{\ell,(r)}=G_t^{\ell,(r)}\otimes A_t^{\ell,(r)}$，其中 A/G 已完成式~\eqref{eq:shrink} 的尺度处理与收缩。固定同一批统计图像，在相邻状态上重测，图中纵轴为
\begin{equation}
 d_r=\mathop{\mathrm{median}}_{t,\ell}
 \left[1-\cos_F\!\left(M_t^{\ell,(r-1)},M_t^{\ell,(r)}\right)\right],
 \qquad r=1,\ldots,4,
 \label{eq:metric-change}
\end{equation}
其中 $\cos_F$ 是将矩阵向量化后的余弦相似度，中位数取有效层/任务记录。横轴 $0$--$1$ 表示初始与第一轮状态，其他位置类推。数值越小，度量方向越接近；它不衡量整体幅度、参数距离或预测误差。这里的 $M$ 来自 KL 修正代理，不是普通 Fisher 矩阵。Kronecker 度量的余弦可由两侧因子的余弦相乘得到，无需构造大矩阵。曲线变小只支持统计趋稳，重测的价值仍由沿用初始统计和同状态干预判断。

\paragraph{多起点适配及其边界。}
表~\ref{tab:starts} 保留全部四种起点。未拉平对照保持聚合任务增量的奇异值比例，逐原始二维矩阵匹配 Iso-C 的增量 Frobenius 范数；控制覆盖 49 个矩阵，包括拼接 Q/K/V，其余参数沿用 Iso-C。两者终态仍有 $8.25$ 个百分点的测试差距，说明总增量能量不能独自解释可修复性。Average/RegMean 的参数距离采用逐层差异除以预训练层权重范数，再取中位数；其 $2.53\%$ 不能解释为全模型欧氏距离。

\begin{table}[htbp]
\centering\small
\caption{完整四起点对照（\%）。诊断与完整测试属于不同分割；四组均由 $\mathcal L_{\mathcal H}$ 选中第四轮。}
\label{tab:starts}
\begin{tabular}{lrrr}
\toprule
初始化 & 初始诊断 & 第四轮诊断 & 终态测试 \\
\midrule
Average & 61.28 & 89.50 & 85.31 \\
RegMean & 83.64 & 89.65 & 85.55 \\
Iso-C & 87.74 & 92.92 & 87.06 \\
未拉平（范数匹配） & 28.86 & 82.18 & 78.80 \\
\bottomrule
\end{tabular}

\end{table}

\paragraph{继续更新与重测统计分别起什么作用？}
对每种初始化先共享第一轮模型 $\theta^{(1)}$，再分成两支。一支沿用初始 $A,G$ 及收缩配置，仍继续求解和选步；另一支在各自当前模型上重新估计。两支复用同一批图像。只解一次从原初始化出发，仅用初始统计做一次共轭梯度（内层 400 步）。表~\ref{tab:refresh} 显示，Average 第一轮之后的 $1.6004$ 个百分点收益中，沿用初始统计也取得 $0.8069$ 个百分点，每轮重测的额外收益为 $0.7935$ 个百分点。

\begin{table}[htbp]
\centering\small
\caption{完整测试准确率（\%）。最后一列为百分点差；只解一次不是等耗时或已收敛对照。}
\label{tab:refresh}
\begin{tabular}{lrrrrr}
\toprule
初始化 & 第一轮 & 四轮沿用初始 & 四轮每轮重测 & 只解一次 & 差值 \\
\midrule
Average & 83.712 & 84.519 & 85.312 & 84.565 & 0.794 \\
RegMean & 85.385 & 85.546 & 85.549 & 85.574 & 0.003 \\
Iso-C & 86.972 & 86.941 & 87.056 & 86.624 & 0.115 \\
\bottomrule
\end{tabular}

\end{table}

\paragraph{统计、方向与接受规则的作用需要区分。}
正文同状态干预保持待更新 checkpoint 一致，但其终态收益包含重测、收缩、有限 CG 和选步的共同作用。将 Average 的\emph{原始}新旧提案匹配中心化 logit RMS 后，$\theta^{(1)}$ 上旧/新方向分别降低独立诊断 KL $0.003379/0.002503$，$\theta^{(2)}$ 上则为 $0.000481/0.002089$。因此，当前统计的接受终态更好，不自动意味着原始方向在每个等幅比较中都更好。现有沿用初始统计的对照也没有单独检验第二轮之后是否每轮都需要重测。

\paragraph{周围网络改变，本块也可能需要适配。}
固定目标 block 的全部权重为初始值，其余网络换成第一轮状态。所有比较从同一个混合模型出发，仅允许目标 block 的线性层更新；分别使用初始统计和在混合模型上重测的统计。Average 的 block 11 自身权重未变，输入与修正梯度余弦却降至 $0.8449$、$0.4915$，使用当前统计也获得更大的独立诊断 KL 降低量。图~\ref{fig:context} 同时保留 block 0 的负结果，支持有条件的网络状态适配。

\begin{figure}[htbp]
\centering
\includegraphics[width=.78\linewidth]{context.pdf}
\caption{\textbf{固定本块权重的跨 block 干预。}Average，seed 0。每个 block 的两根柱共享同一混合起点，纵轴为接受更新后的独立诊断 KL 降低量，越大越好；不同 block 的起点各不相同。三个 block 为预先选定的 0、5、11。}
\label{fig:context}
\end{figure}

\paragraph{归一化的直接消融。}
另一组同机 RTX 4090、seed 0 的配对记录中，归一化与未归一化版本的资产指纹、配置摘要及初始化相同，均迭代四轮并选中第 4 轮；平均测试准确率分别为 $87.0451\%$、$86.6652\%$，最终 $\mathcal L_{\mathcal H}$ 为 $0.155523$、$0.169934$。该消融支持组件有效，但不单独识别卡方相对其他归一化的优势。它与多起点研究的 $87.0558\%$ 属于不同运行记录，不拼接为同一轨迹。上述结果均未提供跨 seed 显著性结论。

\paragraph{跨设定冻结配置。}
表~\ref{tab:protocol} 汇总 CLIP、T5 与 LLM 的校准、教师近似、线搜索与不合并模块。共同流程相同，具体教师与搜索粒度按设定适配。

\begin{table}[htbp]
\centering\small
\caption{跨设定冻结配置。各列是同一迭代框架在不同骨干上的适配，不是三次独立方法。}
\label{tab:protocol}
\setlength{\tabcolsep}{3.2pt}
\renewcommand{\arraystretch}{1.12}
\resizebox{\textwidth}{!}{\begin{tabular}{@{}llll@{}}
\toprule
项 & CLIP ViT-B/32（正式） & FLAN-T5 GLUE & MergeBench LLM \\
\midrule
校准 / 预留 & 每任务 $256/256$ 张 & 每任务 $256/256$ 条 train 提示 & 每域 $256/256$ 条 \texttt{*\_val}，截 1024 token \\
教师 & 专家 + 冻结分类头，逐类精确 & 专家贪心答案的 teacher-forced logits & 专家 top-128 logits \\
起点 & Iso-C，尺度 $1.3$ & Iso-C $1.3$（或已合并模型） & 任意已合并模型；默认 Iso-C $1.3$ \\
求解 & CG 100 步，从当前权重开始 & CG 100 步，从当前权重开始 & CG 100 步，从专家平均权重开始 \\
线搜索 & 全局 + 13 块 & 全局 & 全局 \\
轮数 / 选轮 & 最多 4 轮，$\mathcal L_{\mathcal H}$ 最小 & 同 & 同 \\
不合并 & bias、LayerNorm 继承起点 & 同，嵌入按起点 & \texttt{embed\_tokens}、\texttt{lm\_head} 保留基座 \\
\bottomrule
\end{tabular}
}
\end{table}

\newpage
\section{跨设定主实验补充}
\label{app:cross-setting}

正文表~\ref{tab:ink-vision-summary} 只放覆盖完整的九格规模对照；单设定重复与起点研究另见表~\ref{tab:clip-initializations}。表~\ref{tab:clip-main} 补充固定基准实现的 ViT-B/32 八任务四种子结果与外部参照；表~\ref{tab:clip-scaling} 给出工作台相对各自固定起点的变化、FT 差距及覆盖未齐的外部基线，B/32 八任务的 CM + Iso-C 为 $86.60$。表~\ref{tab:t5-main} 与正文逐任务表共享同一份日志回填数据，原始目标、目标插值 $0.3$ 和 $0.6$ 分行，误差只来自真实重复种子。缺测位置保留空缺，不使用排版占位数。

% Generated by motivation_method/update_results.py from data/verified_results.json.
\begin{table}[!htbp]
\centering\small
\caption{CLIP ViT-B/32 八任务：主配置 CM + Iso-C 与外部参照（\%）。采用固定基准实现，CM 为四种子均值及样本标准差，外部方法为已归档单次结果。}
\label{tab:clip-main}
\begin{tabular}{@{}lcc@{}}
\toprule
方法 & seeds & test 宏平均 \\
\midrule
Individual FT & --- & 89.96 \\
Zero-shot & 1 & $46.69$ \\
WeightAvg & 1 & $66.42$ \\
Task Arith. & 1 & $66.49$ \\
TIES & 1 & $69.41$ \\
Consensus TA & 1 & $70.36$ \\
Iso-C & 1 & $81.33$ \\
Iso-CTS & 1 & $81.04$ \\
TSV-M & 1 & $81.35$ \\
RegMean & 1 & $81.27$ \\
WUDI & 1 & $82.55$ \\
LiNeS (Iso-C) & 1 & $83.37$ \\
CoM & 1 & $83.38$ \\
ESM & 1 & $84.19$ \\
SVC (Iso-CTS) & 1 & $81.61$ \\
SVC (ESM) & 1 & $82.38$ \\
\rowcolor{black!8}
CM + Iso-C & 4 & $87.059\pm 0.021$ \\
Delta vs. FT & --- & -2.90 \\
\bottomrule
\end{tabular}
\end{table}

% Generated by motivation_method/update_results.py from data/verified_results.json.
\begin{table}[!htbp]
\centering\small
\caption{CM 工作台九格（单种子）：固定 Iso-C 起点到留出 KL 选定轮。}
\label{tab:clip-scaling}
\begin{tabular}{@{}lrrr@{}}
\toprule
设定 & 起点 & CM 原型 & 差值 \\
\midrule
ViT-B-32 / 8 & 81.40 & 86.60 & +5.21 \\
ViT-B-32 / 14 & 74.29 & 83.93 & +9.64 \\
ViT-B-32 / 20 & 66.76 & 83.42 & +16.66 \\
ViT-B-16 / 8 & 86.40 & 88.99 & +2.59 \\
ViT-B-16 / 14 & 78.92 & 86.45 & +7.52 \\
ViT-B-16 / 20 & 71.26 & 85.83 & +14.57 \\
ViT-L-14 / 8 & 90.14 & 91.52 & +1.38 \\
ViT-L-14 / 14 & 84.96 & 88.75 & +3.79 \\
ViT-L-14 / 20 & 82.41 & 88.96 & +6.55 \\
\bottomrule
\end{tabular}
\end{table}

% Generated by motivation_method/update_results.py from data/verified_results.json.
\begin{table}[!htbp]
\centering\small
\caption{Flan-T5 GLUE 八任务实测宏平均。原版与目标插值扩展分行；所有误差为样本标准差。large 各行均为单种子。}
\label{tab:t5-main}
\begin{tabular}{@{}lccc@{}}
\toprule
方法 & base & base seeds & large \\
\midrule
Weight Avg. & $78.88$ & 1 & $86.53$ \\
Task Arith. & $78.89$ & 1 & $87.30$ \\
TIES & $80.01$ & 1 & $86.16$ \\
Iso-C 1.3 & $79.64$ & 1 & $82.42$ \\
TSV-M & $82.76$ & 1 & $88.09$ \\
RegMean 0.9 & $83.80$ & 1 & $88.24$ \\
TA + FeatCal（移植） & $84.84\pm 0.22$ & 5 & $89.06$ \\
\rowcolor{black!8}
CM + Iso-C & $84.92\pm 0.05$ & 5 & $88.99$ \\
\rowcolor{black!8}
CM + Iso-C（目标插值 0.3） & $85.64\pm 0.04$ & 4 & $87.97$ \\
\rowcolor{black!8}
CM + Iso-C（目标插值 0.6） & $85.77\pm 0.08$ & 3 & --- \\
\bottomrule
\end{tabular}
\end{table}


\begin{table}[htbp]
\centering\small
\caption{LLM 的起点与后处理结果（单 seed）。同一行比较同一模型及指标；这是起点前后对照，不是所有方法的排名表。}
\label{tab:llm-summary}
\resizebox{\textwidth}{!}{% Generated by motivation_method/update_results.py from data/verified_results.json.
\begin{tabular}{@{}llrrrr@{}}
\toprule
模型 & CM 配置 & \multicolumn{2}{c}{GSM8K} & \multicolumn{2}{c}{IFEval} \\
\cmidrule(lr){3-4}\cmidrule(l){5-6}
 &  & 起点 & + CM & 起点 & + CM \\
\midrule
Llama-3.2-3B & CM + Iso-C 1.3 & 52.84 & 55.65 & 9.80 & 13.12 \\
Llama-3.2-3B & CM + TSV-M & 54.13 & 55.04 & 18.30 & 15.53 \\
Llama-3.2-3B & CM + RegMean 0.5 & 53.53 & 52.39 & 11.46 & 13.12 \\
Gemma-2-2B & CM + Model Soup & 35.25 & 48.07 & 18.11 & 26.06 \\
\bottomrule
\end{tabular}
}
\end{table}

表~\ref{tab:llm-summary} 摘录含收益与下降的行。表~\ref{tab:llama-full} 与表~\ref{tab:gemma-full} 给出 MergeBench 全部已跑格子。Llama 锚点为基座 GSM8K $27.98$、math 专家 $53.22$；三域宏平均是 GSM8K、IFEval 与 multilingual 的等权平均。Iso-C 起点换三个校准 seed 的 GSM8K 为 $55.65/53.98/56.25$，均值 $55.29\pm 1.18$。Gemma 锚点为基座 GSM8K $25.32$ / IFEval $13.68$，math 专家 $54.59$，instruction 专家 IFEval $36.97$。Gemma 上 data-free 的 Localize-and-Stitch 为 $51.33$，高于表中所有后处理结果；同一方法在 Llama 上为 $16.53$。

\begin{table}[htbp]
\centering\small
\caption{Llama-3.2-3B：以官方配置的合并结果为起点，叠加四轮 \method{}（单 seed）。}
\label{tab:llama-full}
\setlength{\tabcolsep}{3.4pt}
\resizebox{\textwidth}{!}{\begin{tabular}{@{}lrrrrrrr@{}}
\toprule
起点（官方超参） & \multicolumn{3}{c}{GSM8K} & \multicolumn{2}{c}{IFEval} & \multicolumn{2}{c}{三域宏平均} \\
\cmidrule(lr){2-4}\cmidrule(lr){5-6}\cmidrule(l){7-8}
 & 起点 & +\,\method{} & $\Delta$ & 起点 & +\,\method{} & 起点 & +\,\method{} \\
\midrule
Iso-C $1.3$（默认） & $52.84$ & $55.65$ & $+2.81$ & $9.80$ & $13.12$ & $36.64$ & $38.63$ \\
TSV-M & $54.13$ & $55.04$ & $+0.91$ & $18.30$ & $15.53$ & $39.97$ & $39.26$ \\
Model Soup & $39.95$ & $54.06$ & $+14.11$ & $9.61$ & $16.64$ & $32.32$ & $39.16$ \\
Task Arithmetic $0.4$ & $40.71$ & $51.86$ & $+11.15$ & $20.70$ & $15.71$ & $36.14$ & $38.20$ \\
TIES $0.3/0.4$ & $42.53$ & $50.27$ & $+7.74$ & $9.80$ & $14.42$ & $33.25$ & $37.10$ \\
RegMean $a{=}0.5$ & $53.53$ & $52.39$ & $-1.14$ & $11.46$ & $13.12$ & $37.51$ & $37.61$ \\
\bottomrule
\end{tabular}
}
\end{table}

\begin{table}[htbp]
\centering\small
\caption{Gemma-2-2B：以官方配置的合并结果为起点，叠加四轮 \method{}（单 seed）。四格 GSM8K 与 IFEval 均为正。}
\label{tab:gemma-full}
\begin{tabular}{@{}lrrrrr@{}}
\toprule
起点（官方超参） & \multicolumn{3}{c}{GSM8K} & \multicolumn{2}{c}{IFEval} \\
\cmidrule(lr){2-4}\cmidrule(lr){5-6}
 & 起点 & +\,\method{} & $\Delta$ & 起点 & +\,\method{} \\
\midrule
Model Soup & $35.25$ & $48.07$ & $+12.82$ & $18.11$ & $26.06$ \\
Task Arithmetic $0.4$ & $39.88$ & $47.46$ & $+7.58$ & $25.32$ & $27.17$ \\
TIES $0.3/0.4$ & $42.91$ & $46.93$ & $+4.02$ & $19.22$ & $26.80$ \\
TSV-M & $41.17$ & $46.32$ & $+5.15$ & $26.25$ & $29.21$ \\
\bottomrule
\end{tabular}

\end{table}

表~\ref{tab:clip-starts} 与表~\ref{tab:t5-starts} 汇总从其它方法输出开始运行 CM 的结果。CLIP 为工作台实现、seed 0、四轮预算，同时报告按 $\mathcal L_{\mathcal H}$ 选轮后的 test 分数和末轮分数；Delta 使用选轮分数与同环境 Individual FT 比较。默认 Iso-C 的选轮分数为 $86.60$，末轮为 $86.65$，均不与固定基准实现的四种子均值 $87.06$ 混用。CM + WeightAvg 为 $85.12$、CM + RegMean 为 $85.32$、CM + TSV-M 为 $86.38$，另保留 Task Arithmetic 与 Zero-shot 起点。

T5 起点对照均为配对 seed 0：base 分别列出原始目标和目标插值 $0.3$，large 列原始目标。base 插值版从 Iso-C 开始为 $85.59$，不能用四种子均值 $85.64$ 替代；原始目标从 FeatCal 和 RegMean 开始分别为 $84.82$、$84.78$。FeatCal 起点本身以 Task Arithmetic 为底座。T5 的已测终点较接近；LLM 仍有起点差异和负向格，不据此声称起点无关。CLIP 上 Iso-C、TSV-M 的选轮终点仍高于其它常用起点。相反顺序的后处理另行区分：先运行 CM 再做 FeatCal，base 原始目标为 $84.64$、插值 $0.3$ 为 $84.78$，large 原始目标为 $89.32$；large 的该结果使用两份校准数据，未进入主表。

\begin{table}[htbp]
\centering\small
\caption{CLIP ViT-B/32 八任务换起点（工作台，同一 seed 0、四轮预算）。留出选轮与末轮均为 test 分数；Delta 相对 Individual FT 为百分点差。}
\label{tab:clip-starts}
% Generated by motivation_method/update_results.py from data/verified_results.json.
\begin{tabular}{@{}lrrrr@{}}
\toprule
配置 & 起点分 & 留出选轮 & 末轮 & Delta vs. FT \\
\midrule
\method{} + Iso-C & 81.40 & 86.60 & 86.65 & -3.36 \\
\method{} + TSV-M & 79.64 & 86.38 & 86.40 & -3.58 \\
\method{} + RegMean & 80.53 & 85.32 & 85.32 & -4.65 \\
\method{} + Zero-shot & 46.70 & 85.25 & 85.25 & -4.71 \\
\method{} + WeightAvg & 62.10 & 85.12 & 85.12 & -4.85 \\
\method{} + Task Arithmetic & 65.54 & 85.09 & 85.09 & -4.88 \\
\bottomrule
\end{tabular}

\end{table}

\begin{table}[htbp]
\centering\small
\caption{Flan-T5 换起点（同一 seed 0）。原始目标与目标插值分列；FeatCal 起点以 Task Arithmetic 为底座。横线表示缺测。}
\label{tab:t5-starts}
\resizebox{\textwidth}{!}{\begin{tabular}{@{}lcccc@{}}
\toprule
起点 & base 起点 & +\,\method{} & large 起点 & +\,\method{} \\
\midrule
Iso-C $1.3$（默认） & $79.64$ & $85.64$ & $82.42$ & $88.99$ \\
TA + FeatCal & $84.56$ & $85.44$ & $89.06$ & $89.10$ \\
RegMean $a{=}0.9$ & $83.80$ & $85.57$ & --- & --- \\
TSV-M & $82.76$ & $85.48$ & $88.09$ & $89.07$ \\
Task Arithmetic $0.3$ & $78.89$ & $85.50$ & --- & --- \\
\bottomrule
\end{tabular}
}
\end{table}

RTX 4090 单卡上，Llama-3.2-3B 每轮约 $35$ 分钟，Gemma-2-2B 约 $20$ 分钟，四轮加 GSM8K 评测约 $2.6$ 小时一格，主机内存峰值约 $85$\,GB。T5-base 每轮约 $1.5$ 分钟，评测约 $6.5$ 分钟。

\newpage
\section{逐层求解：收缩、正规方程与共轭梯度}
\label{app:solver}

正文在当前学生上估计 $A,G$ 后，沿用 MaTS 的 $G\otimes A$ 求解器~\citep{tam2024mats}。与 MaTS 的差别主要在两侧矩阵量的是什么；求解器本身只有细微改动。以下给出尺度处理、层回归代理与有限步 CG。

先统一两侧统计的平均对角元，并仅对输出侧做由样本估计强度的 Ledoit--Wolf 收缩。对非零 $d\times d$ 矩阵 $M$，处理为
\begin{equation}
 \mathcal S_\lambda(M)=(1-\lambda)\frac{M}{\tr(M)/d}+\lambda I,
 \qquad\mathcal S_\lambda(0)=0.
 \label{eq:shrink}
\end{equation}
输入取 $\lambda=0$，输出强度逐任务、逐层估计。以下求解式中的 $A_t^\ell,G_t^\ell$ 均指处理后的矩阵。

当前输入与修正轴给出候选 $X$ 相对专家 $W_t^\ell$ 的平方响应差
$\E[(h^{\ell\top}(X-W_t^\ell)u_t^\ell)^2]$。采用 K-FAC 形式的因子化~\citep{martens2015kfac}，以 $G_t^\ell\otimes A_t^\ell$ 表示双侧权重，构造逐层回归代理
\begin{equation}
 \min_X\ \frac12\sum_{t=1}^T
 \left\|(A_t^\ell)^{1/2}(X-W_t^\ell)(G_t^\ell)^{1/2}\right\|_F^2.
 \label{eq:regression}
\end{equation}
其正规方程沿用相同权重约定：
\begin{equation}
 \sum_{t=1}^T A_t^\ell X G_t^\ell
 =\sum_{t=1}^T A_t^\ell W_t^\ell G_t^\ell.
 \label{eq:normal-equation}
\end{equation}
以当前 $W^{\ell,(r)}$ 为初值，用对角预条件 CG 至多求解 100 步，无需显式构造完整 Kronecker 矩阵。正规方程有形式解，但该算子通常病态，故不用一步求逆，见附录~\ref{app:onestep}。某任务的 $G$ 为零时不参与该层拟合；整层没有有效信号时保留当前权重。

令 $G=I$ 可得到仅含输入权重的 RegMean 类回归形式，但不能据此等同两者的完整实现。式~\eqref{eq:regression} 是结合专家锚点、因子化和统计处理构造的代理，不是网络 KL 的精确二阶展开。

\newpage
\section{为何不把当前正规方程一步求逆}
\label{app:onestep}

\paragraph{闭式修正及其算子。}
附录~\ref{app:solver} 的层回归代理~\eqref{eq:regression} 的正规方程为~\eqref{eq:normal-equation}。令 $L(X)=\sum_{t=1}^T A_t X G_t$。由 $\mathrm{vec}(AXB)=(B^\top\otimes A)\mathrm{vec}(X)$，
\begin{equation}
 \mathrm{vec}\bigl(L(X)\bigr)
 =\Bigl(\sum_{t=1}^T G_t^\top\otimes A_t\Bigr)\mathrm{vec}(X).
 \label{eq:kronecker-sum}
\end{equation}
$G_t$ 为对称外积，故系数矩阵是 $\sum_t G_t\otimes A_t$。若该算子可逆，一步修正即对其求逆。这只是二次代理的形式解：既不保证数值稳定，也不等于网络 KL 的全局最优。

\paragraph{残差几何使算子退化。}
单任务且 $A,G\succ0$ 时，$\kappa(G\otimes A)=\kappa(G)\kappa(A)$。本文的 $G_t=\E[uu^\top]$ 来自尚未消除的 KL 修正统计；师生一致时 $u=0$，有效秩由当前残差轴决定，通常远低于层宽。多任务时 Kronecker 之和一般不能写成单个 Kronecker，显式求逆还要面对 $d_{\rm in}d_{\rm out}$ 维算子。某任务 $G_t=0$ 时该任务本就不进入拟合；接近退化时，不可靠的残差方向对应很小的特征值。

\paragraph{正则化困境。}
为强制可逆而引入 $\lambda I$，或把输出侧写成式~\eqref{eq:shrink} 的 $\mathcal S_\lambda(G)$，会在两个失效端点之间折中。$\lambda\to0$ 时回到近奇异残差 Gram，反演被近零特征值主导，提案沿不可靠修正轴放大。$\lambda\to1$ 时 $\mathcal S_\lambda(G)\to I$，残差几何被冲成各向同性，方程退化为仅由 $A_t$ 加权的输入侧匹配。中间取值只是在放大噪声与抹掉当前需求之间选择，并不能把固定系数的系统变成已经适配后的解。实现中的 Ledoit--Wolf 收缩由样本估计强度决定，目的是让有限 CG 可运行，而不是证明存在一个使一步求逆充分的 $\lambda$。

\paragraph{与有限 CG、选步和重测的关系。}
因此从当前权重出发做对角预条件 CG，截断迭代相当于隐式谱滤波，避免对最小特征值做完全反演；$\mathcal L_{\mathcal H}$ 上的零步长允许拒绝不可靠提案。接受后再测量 $A,G$，使下一轮方程对应新的残差几何。表~\ref{tab:refresh} 的只解一次把同一初始统计解得更充分，Average 上仍低于四轮每轮重测，说明外层适配不能被一步、或更长的内层求解替代。该对照不是条件数测量，也不声称所有层的 $A,G$ 都不可逆。

\newpage
\section{卡方归一化的依据与优势边界}
\label{app:normalization}

以下固定一个样本，设教师 $p$ 为概率分布、学生 $q_c>0$，$e=p-q$，$\mathbf1^\top e=0$，$D=\sum_c e_c^2/q_c$。$F=\mathrm{diag}(q)-qq^\top$，$J$ 为当前学生的 logit Jacobian。所有界均作用于矩阵尺度统一和收缩\emph{之前}的修正外积。

\subsection{学生 Fisher 几何下的精确尺度}
\begin{proposition}[锐利的 Fisher 相对预算]
\label{prop:fisher}
对 $e\ne0$，
\begin{equation}
 D=\sup_{\mathrm{Var}_q(f)>0}
 \frac{(\E_p f-\E_q f)^2}{\mathrm{Var}_q(f)}
 =e^\top F^\dagger e,\qquad
 ee^\top\preceq dF\ \Longleftrightarrow\ d\ge D.
 \label{eq:sharp}
\end{equation}
因此，对 $u=J^\top e/\sqrt{D+\varepsilon}$、$\varepsilon\ge0$，有
\begin{equation}
 uu^\top\preceq\frac{D}{D+\varepsilon}J^\top FJ\preceq J^\top FJ.
\end{equation}
\end{proposition}
\begin{proof}
令 $a_c=p_c/q_c-1$，则 $\E_q a=0$，$\E_q a^2=D$。对任意 $f$，
$\E_p f-\E_q f=\E_q[a(f-\E_q f)]$。
Cauchy--Schwarz 给出平方差至多为 $D\,\mathrm{Var}_q(f)$，取 $f=a$ 达到等号。
又有 $F\,\mathrm{diag}(q)^{-1}e=e$，且 $\ker F=\mathrm{span}(\mathbf1)$。任一解与 $F^\dagger e$ 仅差常数向量，左乘 $e^\top$ 即得 $e^\top F^\dagger e=D$。
矩阵不等式等价于对所有 $f$ 都有 $(f^\top e)^2\le d f^\top Ff$，其最小 $d$ 由刚才的上确界确定。最后除以 $D+\varepsilon$ 并作 $J^\top(\cdot)J$ 合同变换即可。
\end{proof}

这一结论是经典散度变分框架~\citep{nguyen2010divergence} 在当前修正信号上的具体化。它说明根号卡方是保留 logit 残差方向并满足单位学生 Fisher 预算的\emph{最小}分母；原始外积满足同一预算当且仅当 $D\le1$，这里的 1 是预算尺度，非精度阈值。若换成其他分母 $\sqrt{s}$，同一个全方向保证要求 $s\ge D$。例如 $p=(1,0),q=(\delta,1-\delta)$ 时，$D/(2\KL(p\Vert q))\to\infty$，故根号 KL 没有同一全局保证。这不排除它在其他目标上更有效。

\subsection{整体尺度统一后，改变的是方向权重}
将一个图像全部 token 的归一化外积相加，记为 $Z_i$。令 $s_i=D_i+\varepsilon$、$a_i=\tr Z_i$、$P_i=Z_i/a_i$，忽略 $a_i=0$ 的项，且假设总迹非零。归一化与未归一化矩阵的单位迹形式分别为
\begin{equation}
 N=\sum_i b_iP_i,\qquad R=\sum_i\frac{b_i s_i}{\bar s}P_i,
 \quad b_i=\frac{a_i}{\sum_j a_j},\quad\bar s=\sum_i b_i s_i.
\end{equation}
直接展开可得
\begin{equation}
 R-N=\frac{\sum_i b_i(s_i-\bar s)(P_i-N)}{\bar s}
 =\frac{\mathrm{Cov}_b(s,P)}{\bar s}.
 \label{eq:covariance}
\end{equation}
实际的平均对角线统一尺度仅再乘输出维度。因此，归一化的作用不止于改变整体梯度大小，而是移除额外的分歧幅度加权；归一化后的层能量 $a_i$ 仍有差异。若所有 $s_i$ 或所有 $P_i$ 相同，两种形状一致。自动收缩也会随统计变化，故式~\eqref{eq:covariance} 只刻画收缩前的直接作用。

\subsection{何时能得到更好的合并解？}
几何预算本身不保证更高准确率。下面给出一个保留式~\eqref{eq:regression} 结构的可解析条件：固定两个专家 $w_1=e_1=(1,0)^\top,w_2=0$、输入度量均为 1，第二任务输出度量为 $I/2$。第一任务的参考需求为保留 $e_1$ 方向；其估计度量中，非目标方向 $e_2$ 在归一化后占比为 $0<\rho<1$。设两类方向的分歧平方幅度比为
\begin{equation}
 k=\frac{s_{\mathrm{nuisance}}}{s_{\mathrm{useful}}}>0,
 \qquad \tau=\frac{\rho k}{1-\rho+\rho k}.
\end{equation}
则第一任务单位迹度量分别为 $\mathrm{diag}(1-\rho,\rho)$ 与 $\mathrm{diag}(1-\tau,\tau)$。

\begin{proposition}[两专家二次合并的风险边界]
\label{prop:advantage}
定义
\begin{equation}
 L_t(w)=\tfrac12(w-e_1)^\top\mathrm{diag}(1-t,t)(w-e_1)
 +\tfrac12 w^\top(I/2)w,
 \qquad w(t)=\arg\min_w L_t(w).
\end{equation}
以\emph{预先指定}的无污染目标 $L_0$ 为参照，精确求解且不作输出收缩时，
\begin{equation}
 L_0(w(\rho))<L_0(w(\tau))\quad\Longleftrightarrow\quad k>1.
\end{equation}
\end{proposition}
\begin{proof}
Hessian 为 $\mathrm{diag}(3/2-t,1/2+t)\succ0$，故
$w(t)=((1-t)/(3/2-t),0)^\top$，$w(0)=(2/3,0)^\top$。
代入同一参考目标，得到超额风险
\begin{equation}
 \mathcal E(t)=L_0(w(t))-L_0(w(0))
 =\frac34\left(\frac{1-t}{3/2-t}-\frac23\right)^2
 =\frac{t^2}{3(3-2t)^2}.
\end{equation}
它在 $(0,1)$ 上严格递增，而
$\tau-\rho=\rho(1-\rho)(k-1)/(1-\rho+\rho k)$，因此结论成立。
\end{proof}

该构造说明，当非目标方向同时具有更大的分歧幅度时，原始外积会进一步提高其权重，而归一化可以降低合并风险。$k=1$ 持平，$k<1$ 时同一目标下未归一化严格更好。归一化也有非零污染误差 $\mathcal E(\rho)$，因此参照目标并未事后定义为归一化解。命题是二次代理中的条件性结果，实际网络中“高分歧是否意味着低价值”需要独立干预检验。

\paragraph{实现与适用范围。}
使用 $\varepsilon=10^{-8}$ 并 detach 分母；实际 $p\to q$ 时修正趋零。数学证明要求正概率与精确算术，float32 下溢或 tiny clamp 触发时不直接沿用精确 Fisher 恒等式。学生自信但错误也可能使归一化压低重要修正：当 $p=(1,0),q=(\delta,1-\delta)$，KL 发散而 $\|(p-q)/\sqrt{D+\varepsilon}\|_2\sim\sqrt{2\delta}\to0$。此外，自动收缩、有限 CG、非线性传播和选步都会影响最终结果；本附录不提供完整网络准确率的普遍支配保证。
